\subsection{\bt{Lifecycle of one library}{一个定义库的生命周期}}
\label{sec:lifecycle-eval}
% 双语版每格是“英文 / 中文”，单栏放不下，改为通栏
\ifbilingual
\begin{table*}[t]
\caption{\bt{Lifecycle of one shared library: correct answers of each wave of \LcWave\ fresh consumers as updates accumulate (no rollback), and what \system's layer did in the step. One run per method, one million sales rows.}{一个共享定义库的生命周期：更新逐步累积（不回滚），每一步 \LcWave\ 个新使用者答对的题数，以及 \system\ 这一层在该步做了什么。每种方法一次运行，100 万行销售数据。}}
\label{tab:lifecycle}
\footnotesize
\setlength{\tabcolsep}{2.5pt}
% generated by tools/lifecycle-stats.py; do not edit
\begin{tabularx}{\linewidth}{@{}lrrrL@{}}
\toprule
\bt{Step}{步骤} & \bt{\shortstack[r]{No\\memory}}{\shortstack[r]{无\\记忆}} & \bt{\shortstack[r]{Example\\retrieval}}{\shortstack[r]{示例\\检索}} & \bt{\shortstack[r]{MAVRA\\(ours)}}{\shortstack[r]{MAVRA\\（本文）}} & \bt{\system\ actions}{\system\ 的动作}\\
\midrule
\bt{Unchanged data}{数据不变} & 6 & 15 & 15 & \bt{--}{--}\\
\bt{Sales appended}{追加销售} & 7 & 15 & 15 & \bt{--}{--}\\
\bt{Return status rows}{退货状态行} & 4 & 12 & 15 & \bt{2 repaired}{修复 2}\\
\bt{Restated sales}{更正保留旧行} & 7 & 12 & 14 & \bt{4 repaired}{修复 4}\\
\bt{Item SCD Type 2}{商品维度 SCD 2} & 7 & 15 & 14 & \bt{2 repaired}{修复 2}\\
\bt{Date keys re-keyed}{日期键重编号} & 7 & 15 & 14 & \bt{2 reinstated, 5 undone}{恢复 2，撤销 5}\\
\bt{Duplicate load}{重复加载} & 7 & 9 & 10 & \bt{4 invalidated}{失效 4}\\
\bt{Duplicates removed}{删除重复批次} & 7 & 9 & 12 & \bt{1 relearned, 2 refused}{重新学习 1，拒绝 2}\\
\bt{In-place correction}{原地更正} & 9 & 9 & 13 & \bt{1 relearned}{重新学习 1}\\
\midrule
\bt{Correct, all steps}{全部步骤正确率} & 45\% & 82\% & 90\% & \\
\bt{k tokens per task}{每题 k token} & 37.5 & 18.0 & 28.3 & \\
\bt{k tokens per correct answer}{每个正确答案 k token} & 85.2 & 23.2 & 36.2 & \\
\bottomrule
\end{tabularx}

\end{table*}
\else
\begin{table}[t]
\caption{\bt{Lifecycle of one shared library: correct answers of each wave of \LcWave\ fresh consumers as updates accumulate (no rollback), and what \system's layer did in the step. One run per method, one million sales rows.}{一个共享定义库的生命周期：更新逐步累积（不回滚），每一步 \LcWave\ 个新使用者答对的题数，以及 \system\ 这一层在该步做了什么。每种方法一次运行，100 万行销售数据。}}
\label{tab:lifecycle}
\footnotesize
\setlength{\tabcolsep}{2.5pt}
% generated by tools/lifecycle-stats.py; do not edit
\begin{tabularx}{\linewidth}{@{}lrrrL@{}}
\toprule
\bt{Step}{步骤} & \bt{\shortstack[r]{No\\memory}}{\shortstack[r]{无\\记忆}} & \bt{\shortstack[r]{Example\\retrieval}}{\shortstack[r]{示例\\检索}} & \bt{\shortstack[r]{MAVRA\\(ours)}}{\shortstack[r]{MAVRA\\（本文）}} & \bt{\system\ actions}{\system\ 的动作}\\
\midrule
\bt{Unchanged data}{数据不变} & 6 & 15 & 15 & \bt{--}{--}\\
\bt{Sales appended}{追加销售} & 7 & 15 & 15 & \bt{--}{--}\\
\bt{Return status rows}{退货状态行} & 4 & 12 & 15 & \bt{2 repaired}{修复 2}\\
\bt{Restated sales}{更正保留旧行} & 7 & 12 & 14 & \bt{4 repaired}{修复 4}\\
\bt{Item SCD Type 2}{商品维度 SCD 2} & 7 & 15 & 14 & \bt{2 repaired}{修复 2}\\
\bt{Date keys re-keyed}{日期键重编号} & 7 & 15 & 14 & \bt{2 reinstated, 5 undone}{恢复 2，撤销 5}\\
\bt{Duplicate load}{重复加载} & 7 & 9 & 10 & \bt{4 invalidated}{失效 4}\\
\bt{Duplicates removed}{删除重复批次} & 7 & 9 & 12 & \bt{1 relearned, 2 refused}{重新学习 1，拒绝 2}\\
\bt{In-place correction}{原地更正} & 9 & 9 & 13 & \bt{1 relearned}{重新学习 1}\\
\midrule
\bt{Correct, all steps}{全部步骤正确率} & 45\% & 82\% & 90\% & \\
\bt{k tokens per task}{每题 k token} & 37.5 & 18.0 & 28.3 & \\
\bt{k tokens per correct answer}{每个正确答案 k token} & 85.2 & 23.2 & 36.2 & \\
\bottomrule
\end{tabularx}

\end{table}
\fi

\bi{\emph{Setup (Q8).} The scenarios of Section~\ref{sec:scenarios} roll back every change. Here one library per method lives through \LcSteps\ steps on one million rows without rollback (Table~\ref{tab:lifecycle}): each step applies its update and then sends a wave of \LcWave\ fresh consumers, \LcConc\ at a time, one per scored question. \system\ runs in its full configuration, with optimization, same-snapshot execution, and transactional versions, and its built-in learner relearns any metric left without a valid definition. \emph{Result.} \system\ answers \LcAccCond\% of the \LcTasks\ questions, example retrieval \LcAccTraj\%, and agents without memory \LcAccNone\%. The layer repaired \LcRepaired\ definitions; when the re-keying broke the premise of the \LcOptimized\ optimized definitions, it reinstated \LcByPred\ predecessors and undid the rule for \LcByUndo\ revisions repaired since, and \LcRekeyCorrectCond\ of \LcWave\ answers stayed correct; it invalidated \LcRevoked\ definitions under the duplicate load and relearned their metrics after the cleanup, refusing \LcRelearnRefused\ faulty extracted candidates on the way. Until the invalidation, \LcPreDeclaredPctCond\% of consumers declared the revision they used, and each needed \LcPreTurnsMinCond--\LcPreTurnsMaxCond\ turns and \LcPreTokMinCond--\LcPreTokMaxCond\,k tokens, against \LcPreTurnsMinTraj--\LcPreTurnsMaxTraj\ turns and \LcPreTokMinTraj--\LcPreTokMaxTraj\,k with retrieved examples. The invalidation sends consumers back to the data at \LcDupTokCond\,k tokens per task (\LcDupTokTraj\,k with retrieval), so retrieval needs fewer tokens per correct answer over the run (\LcTokCorrectTraj\ against \LcTokCorrectCond\,k, building and relearning included; \LcTokCorrectNone\,k without memory), while its stale examples answer \LcAfterCorrectMaxTraj\ of \LcWave\ in both steps after the duplicate load, against \LcAfterCorrectMinCond\ and \LcAfterCorrectMaxCond\ for \system.}
{\emph{设置（Q8）。}第~\ref{sec:scenarios} 节的场景在每次变化后回滚。这里每种方法的一个定义库在 100 万行数据上依次经历 \LcSteps\ 步，不回滚（表~\ref{tab:lifecycle}）：每一步先施加更新，再来一波 \LcWave\ 个新使用者，每道计分题一个，\LcConc\ 个同时进行。\system\ 使用完整配置，包括优化、同快照执行与事务性版本，其内置学习者重新学习没有有效定义的指标。\emph{结果。}\system\ 答对 \LcTasks\ 道题中的 \LcAccCond\%，示例检索 \LcAccTraj\%，无记忆的智能体 \LcAccNone\%。这一层修复了 \LcRepaired\ 个定义；日期键重编号破坏 \LcOptimized\ 个优化定义的前提时，它恢复了 \LcByPred\ 个前一修订，对此后被修复过的 \LcByUndo\ 个修订撤销改写规则，\LcWave\ 个答案中 \LcRekeyCorrectCond\ 个保持正确；重复加载下它使 \LcRevoked\ 个定义失效，并在清理之后重新学习这些指标，途中拒绝了 \LcRelearnRefused\ 个提取有误的候选。失效之前，\LcPreDeclaredPctCond\% 的使用者声明了所用的修订，每题用 \LcPreTurnsMinCond--\LcPreTurnsMaxCond\ 轮、\LcPreTokMinCond--\LcPreTokMaxCond\,k token，检索示例时为 \LcPreTurnsMinTraj--\LcPreTurnsMaxTraj\ 轮、\LcPreTokMinTraj--\LcPreTokMaxTraj\,k。失效把使用者送回数据，每题 \LcDupTokCond\,k token（检索为 \LcDupTokTraj\,k），因此在整次运行上检索每个正确答案用的 token 更少（\LcTokCorrectTraj\ 对 \LcTokCorrectCond\,k，含建立与重新学习；无记忆为 \LcTokCorrectNone\,k），而它过期的示例在重复加载之后的两步中都只答对 \LcWave\ 道中的 \LcAfterCorrectMaxTraj\ 道，\system\ 为 \LcAfterCorrectMinCond\ 道和 \LcAfterCorrectMaxCond\ 道。}
